% !TEX root = ../../main.tex

\renewcommand{\thesection}{1.6}

\section{重因式}\label{sec:1.6}

\begin{ex}{1}\label{ex:1.6.1}
    求下列多项式在有理数域上的标准分解:
    \[
    x^5-6x^4+16x^3-24x^2+20x-8.
    \]
\end{ex}

\begin{solution}
    设 \(f(x)=x^5-6x^4+16x^3-24x^2+20x-8\), 则 \(f'(x)=5x^4-24x^3+48x^2-48x+20\).

    通过辗转相除法求得
    \[
    (f(x),f'(x))=x^2-2x+2.
    \]
    故 \(f(x)\) 有重因式 \(x^2-2x+2\), 即
    \[
    (x^2-2x+2)^2\mid f(x).
    \]
    通过长除法求得
    \[
    \frac{f(x)}{(x^2-2x+2)^2}=x-2.
    \]
    故
    \[
    f(x)=(x^2-2x+2)^2(x-2).\qedhere
    \]
\end{solution}

\begin{ex}{8}\label{ex:1.6.8}
    设 \(f(x)\) 是 \(n\) 次多项式, \(n\geq1\).
    证明 \(f'(x)\mid f(x)\) 当且仅当 \(f(x)=a(x-b)^n\),
    其中 \(a,b\) 是数且 \(a\neq0\).
\end{ex}

\begin{proof}
    充分性. \(f'(x)=an(x-b)^{n-1}\).
    因此
    \[
    (x-b)f'(x)=nf(x),
    \]
    故有
    \[
    f'(x)\mid nf(x).
    \]
    由 \(n\) 为正整数, 有 \((f'(x),n)=1\), 故
    \[
    f'(x)\mid f(x).
    \]
    必要性. 设 \(f_1\) 为与 \(f\) 相伴的首一多项式. 设 \(f_1\) 的标准分解为
    \[
    f_1=\prod_{i=1}^{t}p_i^{k_i},
    \]
    则
    \[
    (f_1,f_1')=\prod_{i=1}^{t}p_i^{k_i-1},
    \]
    故
    \begin{equation}
        \frac{f_1}{(f_1,f_1')}=\prod_{i=1}^{t}p_i.
        \label{eq:1.6.8-Multiplication}
    \end{equation}
    设 \(f=af_1\), 则 \(f'=af_1'\), 故有 \(f_1'\) 的首项系数为 \(n\), 且
    \[
    f_1'\mid f_1,
    \]
    故\((f_1,f_1')=\frac{1}{n}f_1'\). 设 \(f_1=gf_1'\),
    则 \(g\) 的首项系数为 \(\frac{1}{n}\), 且
    \[
    \deg g=\deg f_1-\deg f_1'=\deg f-\deg f'=n-(n-1)=1.
    \]
    设 \(g(x)=\frac{1}{n}(x-b)\), 由式~\eqref{eq:1.6.8-Multiplication}, 有
    \[
    \prod_{i=1}^{t}p_i(x)=ng(x)=x-b.
    \]
    故 \(t=1,p_1(x)=x-b,f_1(x)=(x-b)^n,f(x)=a(x-b)^n\).
\end{proof}

\begin{ex}{10}\label{ex:1.6.10}
    设
    \[
    f(x)=\sum_{i=0}^n\frac{x^i}{i!},
    \]
    证明 \(f(x)\) 在任何数域上均无重因式.
\end{ex}

\begin{proof}
    \(f(x)\) 的导式
    \[
    f'(x)=\sum_{i=1}^n\frac{x^{i-1}}{(i-1)!}=\sum_{i=0}^{n-1}\frac{x^i}{i!}.
    \]
    则
    \[
    (f(x),f'(x))=(f(x),f(x)-f'(x))=\p{f(x),\frac{x^n}{n!}}=(f(x),x^n).
    \]
    由于 \(f(0)=1\neq0\), 故 \((f(x),x)=1\), 故 \((f(x),x^n)=1\),
    故 \((f(x),f'(x))=1\), 因此 \(f(x)\) 在任何数域上均无重因式.
\end{proof}